% Charge puis décharge d'un condensateur dans un circuit RC (E = 5 V) (L'essentiel, chapitre 26)
\begin{tikzpicture}
\begin{axis}[otMesures,otAxesBout,otLecture,width=110mm,height=50mm,xmin=0,xmax=10.6,ymin=0,ymax=5.8,
  xtick={0,1,...,10},ytick={0,1,...,5},
  xlabel={$t/\tau$},ylabel={$u_C$ (V)}]
\addplot[otMuted,line width=0.5pt,dash pattern=on 2pt off 1.5pt,domain=0:5,samples=2] {5};
\addplot[otNavy,line width=1pt,domain=0:5,samples=80] {5*(1-exp(-x))};
\addplot[otRed!80!black,line width=1pt,domain=5:10.5,samples=80] {4.97*exp(-(x-5))};
\addplot[otMuted,line width=0.6pt,dash pattern=on 2pt off 1.5pt,domain=0:1,samples=2] {5*x};
\draw[otGreenDark,line width=0.5pt,dash pattern=on 1.5pt off 1pt] (axis cs:1,0) -- (axis cs:1,3.16) -- (axis cs:0,3.16);
\fill[otGreenDark] (axis cs:1,3.16) circle (1.5pt);
\node[pcLegende,anchor=west,text=otGreenDark] at (axis cs:1.1,2.9) {63~\% de $E$ à $t=\tau$};
\draw[otGreenDark,line width=0.5pt,dash pattern=on 1.5pt off 1pt] (axis cs:6,0) -- (axis cs:6,1.83);
\fill[otGreenDark] (axis cs:6,1.83) circle (1.5pt);
\node[pcLegende,anchor=west,text=otGreenDark] at (axis cs:6.1,2.1) {37~\% de $E$ après une durée $\tau$};
\node[pcLegende,anchor=south,text=otNavy] at (axis cs:2.6,5.05) {charge (position 1)};
\node[pcLegende,anchor=south,text=otRed!80!black] at (axis cs:7.9,3.6) {décharge (position 2)};
\end{axis}
\end{tikzpicture}
\par\smallskip\legende{La tension évolue exponentiellement ; la tangente à l'origine coupe l'asymptote à $t=\tau=RC$, et au bout de $5\tau$ le régime permanent est pratiquement atteint.}
